\documentclass[10pt,a4paper]{amsart}
\usepackage[margin=19mm]{geometry}
\usepackage{amsmath,amssymb,mathtools}
\usepackage[scr=rsfs,cal=euler]{mathalfa}
\usepackage{xfrac}
\usepackage[unicode,hidelinks,bookmarks=false]{hyperref}
\newcommand{\ie}{i.e.\ }
\newcommand{\cf}{cf.\ }
\newcommand{\resp}{respectively\ }
\newcommand{\Subsection}{\subsection}
\providecommand{\nobreakdash}{\nobreak\hbox{-}}
\setlength{\emergencystretch}{2em}
\multlinegap0pt
\usepackage{amssymb}
\usepackage{dsfont}
\usepackage{algorithm}\usepackage{algpseudocode}
\usepackage[shortlabels]{enumitem}
\usepackage{xcolor}
\usepackage{tikz}
\newtheorem{lemma}{Lemma}
\newtheorem{proposition}{Proposition}
\title{Packing arithmetic progressions}
\author{Noga Alon, Micha{\l} D\k ebski, Jaros\l aw Grytczuk, Jakub Przyby\l o}
\date{Source: arXiv:2603.02786v1 (CC BY 4.0)}
\begin{document}
\maketitle

\noindent\emph{Source and licence.} Packing arithmetic progressions. Version arXiv:2603.02786v1 (CC BY 4.0), distributed by arXiv under the Creative Commons Attribution 4.0 International licence, http://creativecommons.org/licenses/by/4.0/, as recorded in the arXiv licence metadata for this exact version and shown on the abstract page https://arxiv.org/abs/2603.02786v1. Full licence notice: \texttt{licenses/arxiv-2603.02786-CC-BY-4.0.txt}.\par\smallskip

\noindent\textit{Editorial scope.} Counted unit: the article's proposition p32 (source0.tex lines 613--617). The article's complete original proof of it is delivered with it (source0.tex lines 618--636, one \texttt{proof} environment of 19 lines). No mathematics is written, repaired or re-derived: the statement and the proof are the article's own lines, in the article's own notation, and no retained line is substituted or re-flowed.  Retained uncounted as the source-local context the proof needs: lines 206--209.  Nothing was omitted from any retained span, so the package carries no omission record.  No retained line contains a citation, so the wrapper carries no bibliography and no bibliography entry.  No retained line contains a figure environment, an image-inclusion command or drawing code, so no image is delivered and the package declares no asset dependency.  Editorial formatting only, in addition: an AMS \texttt{amsart} preamble that restates the article's own declarations and macros used by the retained lines, namely the environments \texttt{lemma}, \texttt{proposition} on separate counters, together with the standard packages those lines need.  The retained environments print in this document as Lemma 1, Proposition 1 in order of appearance, whereas the article prints the same items with the numbering of its own full document.  The complete native source of the article as distributed in the arXiv e-print for this exact version is bundled unchanged as \texttt{sources/195638/source0.tex}.
\begin{lemma}
	\label{l23}
	Let $A$ and $B$ be two integer subsets of an interval $I$ of length $g>n$, and assume that $B$ fits in an interval of length $n$. Suppose, further, that $|A| \cdot |B| < g-n$. Then there is a shifted 	copy $z+B$ of $B$ which is disjoint from $A$ and fits in $I$.
\end{lemma}

\begin{proposition}
	\label{p32}
	Let $C_i$, $1 \leqslant i \leqslant n$, be an arbitrary collection of subsets of $[n]$, where for all admissible $i$, $|C_i|=\lfloor \frac{n}{i} \rfloor$. Then there are shifts $r_i+C_i$ that can be packed in an interval
	of length at most $(1+1/e+o(1))n^{3/2}$. This is tight up to the constant factor, that is, there are such sets $C_i$ that cannot be packed in an interval of length smaller than $(1+o(1))n^{3/2}/2$.
\end{proposition}

\begin{proof}
	The first $\sqrt n$ sets $C_i$ can clearly be packed in an interval of length $n^{3/2}$ by shifting each $C_i$ to a separate interval of length $n$. Starting  then with an empty new interval of length 
	$bn^{3/2}$ we pack shifts of the sets $C_i$ for $i > \sqrt n$ one by one, in order, using Lemma \ref{l23}. When trying to embed the set $C_x$, the total size of the sets embedded so far (in this new interval) is
	at most 
	$$
	\sum_{\sqrt n < i<x} \frac{n}{i} =
	(1+o(1)) n (\ln x - \ln {\sqrt n}).
	$$
	(We assume here that, say, $x \geqslant 1.1 \sqrt n $, it is easy to see that there is no problem to embed the earlier sets.)
	
	By Lemma \ref{l23} the embedding can be performed provided for every relevant $x$,
	\begin{equation}
		\label{e31}
		(1+o(1)) n (\ln x - \ln {\sqrt n}) \cdot \frac{n}{x} < b n^{3/2}-n.
	\end{equation}
	The function $f(x)=\frac{\ln x -\ln \sqrt n}{x}$ attains its maximum for $x > \sqrt n$ at $x=e \sqrt n$, where its value is $1/(e \sqrt n)$. Therefore, the inequality in (\ref{e31}) holds for any fixed $b>1/e$ (and sufficiently large $n$). This establishes the first part of the proposition.
	
	To see that it is tight up to a constant factor consider an example in which each of the first $\sqrt n/2$ sets $C_i$ contains the interval	$[\sqrt n]=\{1,2, \ldots ,\sqrt n\}$ as well as the arithmetic progression $\{\sqrt n, 2 \sqrt n, 3 \sqrt n, \ldots,n\}$ (assume, for simplicity, that $\sqrt n $ is an integer.) Then in any disjoint packing of shifts of these $\sqrt n/2$ sets $C_i$, no two intervals of length $n$ containing the shifted $C_i$ can have a nonempty overlap. This completes the proof of the proposition.
\end{proof}

\end{document}
